% =====================================================================
%  Fiche de revision -- Databricks Certified Data Engineer Associate
%  Basee sur le guide officiel (version du 4 mai 2026) + AI Prep Guide
%  Compilation : pdflatex revision-databricks-de-associate.tex  (x2)
% =====================================================================
\documentclass[11pt,a4paper]{article}

% ---- Encodage & langue -------------------------------------------------
\usepackage[utf8]{inputenc}
\usepackage[T1]{fontenc}
\usepackage[french]{babel}
\usepackage{lmodern}

% ---- Mise en page ------------------------------------------------------
\usepackage[a4paper,margin=2.2cm]{geometry}
\usepackage{titlesec}
\usepackage{enumitem}
\usepackage{booktabs}
\usepackage{array}
\usepackage{longtable}
\usepackage{amssymb}
\usepackage{xcolor}
\usepackage{tcolorbox}
\usepackage{listings}
\usepackage{fancyhdr}
\usepackage{hyperref}
\tcbuselibrary{skins,breakable}

\hypersetup{
  colorlinks=true,
  linkcolor=databricksred,
  urlcolor=databricksred,
  citecolor=databricksred,
  pdftitle={Revision Databricks Certified Data Engineer Associate},
  pdfauthor={Fiche de revision}
}

% ---- Couleurs ----------------------------------------------------------
\definecolor{databricksred}{RGB}{255,54,33}
\definecolor{darkgray}{RGB}{51,51,51}
\definecolor{lightgray}{RGB}{245,245,245}
\definecolor{codebg}{RGB}{248,248,248}
\definecolor{kwcolor}{RGB}{0,102,153}
\definecolor{strcolor}{RGB}{163,21,21}
\definecolor{comcolor}{RGB}{0,128,0}
\definecolor{piegebg}{RGB}{255,244,229}
\definecolor{piegeframe}{RGB}{230,126,34}
\definecolor{clebg}{RGB}{232,245,233}
\definecolor{cleframe}{RGB}{46,125,50}
\definecolor{qcmbg}{RGB}{232,240,254}
\definecolor{qcmframe}{RGB}{25,103,210}

% ---- En-tetes / pieds de page -----------------------------------------
\setlength{\headheight}{24pt}
\addtolength{\topmargin}{-12pt}
\pagestyle{fancy}
\fancyhf{}
\fancyhead[L]{\small\textcolor{databricksred}{Databricks Data Engineer Associate}}
\fancyhead[R]{\small\thepage}
\fancyfoot[C]{\small\textit{Fiche de revision -- version examen 4 mai 2026}}
\renewcommand{\headrulewidth}{0.4pt}
\renewcommand{\headrule}{\hbox to\headwidth{\color{databricksred}\leaders\hrule height \headrulewidth\hfill}}

% ---- Style des titres --------------------------------------------------
\titleformat{\section}
  {\normalfont\Large\bfseries\color{databricksred}}
  {\thesection.}{0.6em}{}
\titleformat{\subsection}
  {\normalfont\large\bfseries\color{darkgray}}
  {\thesubsection}{0.6em}{}

% ---- Style du code -----------------------------------------------------
\lstdefinestyle{pystyle}{
  language=Python,
  backgroundcolor=\color{codebg},
  basicstyle=\ttfamily\small,
  keywordstyle=\color{kwcolor}\bfseries,
  stringstyle=\color{strcolor},
  commentstyle=\color{comcolor}\itshape,
  numbers=left, numberstyle=\tiny\color{gray}, numbersep=8pt,
  showstringspaces=false, breaklines=true, frame=single,
  rulecolor=\color{gray!40}, tabsize=2, captionpos=b,
  morekeywords={spark,readStream,writeStream,DataFrame,cloudFiles,option,format,load,save,toTable}
}
\lstdefinestyle{sqlstyle}{
  language=SQL,
  backgroundcolor=\color{codebg},
  basicstyle=\ttfamily\small,
  keywordstyle=\color{kwcolor}\bfseries,
  stringstyle=\color{strcolor},
  commentstyle=\color{comcolor}\itshape,
  numbers=left, numberstyle=\tiny\color{gray}, numbersep=8pt,
  showstringspaces=false, breaklines=true, frame=single,
  rulecolor=\color{gray!40}, tabsize=2, captionpos=b,
  morekeywords={COPY,INTO,MERGE,MATCHED,GRANT,REVOKE,DENY,CATALOG,SCHEMA,VOLUME,STREAMING,MATERIALIZED,CLUSTER,LIQUID,PREDICTIVE,OPTIMIZATION,EXTERNAL,MANAGED,MASK,ROW,FILTER}
}
\lstset{style=pystyle}

% ---- Environnements encadres ------------------------------------------
\newtcolorbox{piege}[1][]{
  breakable, colback=piegebg, colframe=piegeframe,
  fonttitle=\bfseries, title={\faWarning Piege d'examen #1},
  boxrule=1pt, arc=2mm, left=3mm, right=3mm, top=2mm, bottom=2mm
}
\newtcolorbox{cle}[1][]{
  breakable, colback=clebg, colframe=cleframe,
  fonttitle=\bfseries, title={A retenir #1},
  boxrule=1pt, arc=2mm, left=3mm, right=3mm, top=2mm, bottom=2mm
}
\newtcolorbox{qcm}[1][]{
  breakable, colback=qcmbg, colframe=qcmframe,
  fonttitle=\bfseries, title={QCM #1},
  boxrule=1pt, arc=2mm, left=3mm, right=3mm, top=2mm, bottom=2mm
}

% faWarning fallback (pas de fontawesome requis)
\providecommand{\faWarning}{\textbf{!} }

% =====================================================================
\begin{document}

% ---------------------------------------------------------------- GARDE
\begin{titlepage}
\centering
\vspace*{2cm}
{\color{databricksred}\rule{\textwidth}{2pt}}\\[0.5cm]
{\Huge\bfseries Databricks Certified\\[0.3cm] Data Engineer Associate}\\[0.4cm]
{\LARGE Fiche de revision complete}\\[0.3cm]
{\color{databricksred}\rule{\textwidth}{2pt}}\\[1.2cm]

{\large Version de l'examen : \textbf{4 mai 2026}}\\[0.3cm]
{\large Basee sur le guide officiel Databricks + AI Prep Guide}\\[2cm]

\begin{tcolorbox}[width=0.8\textwidth,colback=lightgray,colframe=databricksred,
  title=\textbf{En bref},fonttitle=\large]
\begin{tabular}{@{}ll@{}}
\textbf{Questions notees} & 45 QCM \\
\textbf{Duree} & 90 minutes \\
\textbf{Prix} & 200 USD (+ taxes) \\
\textbf{Format} & En ligne ou centre de test \\
\textbf{Aides} & Aucune autorisee \\
\textbf{Prerequis} & Aucun (6 mois de pratique recommandes) \\
\textbf{Validite} & 2 ans (recertification requise) \\
\end{tabular}
\end{tcolorbox}

\vfill
{\small Document de revision personnel -- verifier toujours sur \url{https://docs.databricks.com}}
\end{titlepage}

% ------------------------------------------------------------ SOMMAIRE
\tableofcontents
\newpage

% ============================================ PONDERATIONS
\section*{Carte des priorites : ponderation par domaine}
\addcontentsline{toc}{section}{Carte des priorites}

Le temps de revision doit suivre le poids de chaque domaine. \textbf{Transformation (22\%)
et Ingestion (21\%) representent a eux seuls 43\% de l'examen} : c'est la ou il faut
investir le plus.

\begin{center}
\renewcommand{\arraystretch}{1.3}
\begin{tabular}{@{}p{9cm}>{\centering\arraybackslash}p{2.5cm}>{\centering\arraybackslash}p{2.5cm}@{}}
\toprule
\textbf{Domaine} & \textbf{Poids} & \textbf{$\approx$ questions} \\
\midrule
\S3 -- Data Transformation and Modeling & \textbf{22\%} & $\sim$10 \\
\S2 -- Data Ingestion and Loading       & \textbf{21\%} & $\sim$9 \\
\S4 -- Working with Lakeflow Jobs        & 16\% & $\sim$7 \\
\S7 -- Governance and Security           & 15\% & $\sim$7 \\
\S5 -- Implementing CI/CD                & 10\% & $\sim$5 \\
\S6 -- Troubleshooting, Monitoring, Optim.& 10\% & $\sim$5 \\
\S1 -- Databricks Intelligence Platform  & 6\%  & $\sim$3 \\
\midrule
\textbf{Total} & \textbf{100\%} & \textbf{45} \\
\bottomrule
\end{tabular}
\end{center}

% ============================================ RENOMMAGES
\section{Pieges de renommage (a lire en premier)}

Les modeles d'IA ont ete entraines \emph{avant} les changements recents de la plateforme
Databricks : ils utilisent souvent les anciens noms. \textbf{L'examen utilise les noms
actuels}. Memorise ce tableau -- c'est la cause n\textsuperscript{o}1 d'echec en preparation avec l'IA.

\begin{center}
\renewcommand{\arraystretch}{1.4}
\begin{longtable}{@{}p{4.3cm}p{4.3cm}p{5.5cm}@{}}
\toprule
\textbf{Ancien nom} & \textbf{Nom actuel (2026)} & \textbf{Role} \\
\midrule
\endhead
Databricks Workflows & \textbf{Lakeflow Jobs} & Orchestration de pipelines (DAG de taches) \\
Delta Live Tables (DLT) & \textbf{Lakeflow Spark Declarative Pipelines} & Pipelines declaratifs ETL \\
Databricks Repos & \textbf{Git Folders} & Integration Git dans le workspace \\
Databricks Asset Bundles (DAB) & \textbf{Automation Bundles} (Declarative) & Packaging / promotion dev-test-prod \\
Partner Connect / connecteurs & \textbf{Lakeflow Connect} & Ingestion managee depuis sources externes \\
\bottomrule
\end{longtable}
\end{center}

\begin{cle}
Si une reponse de QCM emploie un ancien nom (Workflows, DLT, Repos) alors qu'une autre
emploie le nom actuel pour la meme chose, l'examen tend a privilegier la terminologie
\textbf{a jour}. Mais attention : le concept prime sur le nom -- le fond ne change pas.
\end{cle}

% ============================================ S1 PLATFORM
\section{Databricks Intelligence Platform (6\%)}

\subsection{Architecture en couches}

La plateforme separe \textbf{control plane} (gere par Databricks : UI, orchestration,
metadonnees) et \textbf{compute/data plane} (dans ton cloud : clusters, stockage objet
ADLS/S3/GCS).

\subsection{Delta Lake}

Format de table par defaut. Points cles pour l'examen :
\begin{itemize}[leftmargin=*,nosep]
  \item \textbf{Transactions ACID} sur le stockage objet (via le transaction log \texttt{\_delta\_log}).
  \item \textbf{Time travel} : interroger une version anterieure (\texttt{VERSION AS OF}, \texttt{TIMESTAMP AS OF}) $\rightarrow$ rollback fiable apres un mauvais ingest.
  \item \textbf{Schema enforcement} + \textbf{schema evolution}.
  \item \textbf{MERGE} (upsert), \texttt{OPTIMIZE}, \texttt{VACUUM}.
\end{itemize}

\subsection{Unity Catalog (UC)}

Couche de gouvernance centralisee. Hierarchie a 3 niveaux :
\[
\texttt{catalog} \;\rightarrow\; \texttt{schema (database)} \;\rightarrow\; \texttt{table / view / volume / model}
\]
UC fournit : controle d'acces fin, lineage, audit, partage. \textbf{Source unique de verite}
pour l'IA et le BI.

\subsection{Services de compute (clusters)}

\begin{center}
\renewcommand{\arraystretch}{1.3}
\begin{tabular}{@{}p{3.5cm}p{11cm}@{}}
\toprule
\textbf{Type} & \textbf{Usage / caracteristiques} \\
\midrule
All-purpose & Developpement interactif, notebooks partages, exploration. Plus cher. \\
Job cluster & Cree pour un job, detruit a la fin. Moins cher pour l'ETL planifie. \\
SQL Warehouse & Requetes SQL/BI, demarrage rapide, multi-utilisateurs concurrents. \\
Single-node & Petits tests, dev leger, pas de workers. \\
\bottomrule
\end{tabular}
\end{center}

\begin{piege}
Pour \textbf{plusieurs analystes lancant des requetes SQL ad hoc} avec demarrage rapide et
concurrence : la bonne reponse penche vers un \textbf{SQL Warehouse} (ou cluster high-concurrency
avec autoscaling) -- \emph{pas} un job cluster (fait pour l'ETL planifie) ni un single-node.
\end{piege}

\begin{qcm}
Un ingenieur veut : iteration rapide, rollback fiable apres mauvais ingest, piste d'audit
reglementaire, et source unique de verite pour IA + BI. Quelle strategie ?

\smallskip
A. Stockage CSV sur DBFS + copies nocturnes \hfill
B. \textbf{Delta Lake (ACID + time travel) gouverne par Unity Catalog} \\
C. Stockage objet seul + requetes SQL ad hoc \hfill
D. DataFrames en memoire ephemeres

\smallskip
\textit{Reponse : \textbf{B}. Delta = ACID + time travel (rollback, audit) ; UC = acces
coherent et lineage. Les autres options n'offrent ni transactions ni gouvernance.}
\end{qcm}

% ============================================ S2 INGESTION
\section{Data Ingestion and Loading (21\%)}

\subsection{Les patterns d'ingestion}

Batch / streaming / incremental. Le choix depend de : volume, frequence, types de donnees,
besoins de gouvernance.

\subsection{COPY INTO}

Commande SQL \textbf{idempotente} : charge de maniere incrementale des fichiers depuis le
stockage objet vers une table UC. Retient les fichiers deja charges $\rightarrow$ pas de doublons.

\begin{lstlisting}[style=sqlstyle,caption={COPY INTO depuis le stockage objet}]
COPY INTO mon_catalogue.bronze.ventes
FROM 's3://bucket/ventes/'
FILEFORMAT = CSV
FORMAT_OPTIONS ('header' = 'true', 'inferSchema' = 'true')
COPY_OPTIONS ('mergeSchema' = 'true');
\end{lstlisting}

\subsection{Auto Loader (\texttt{cloudFiles})}

Ingestion \textbf{incrementale en streaming} de nouveaux fichiers. Gere schema inference,
schema evolution et checkpointing. Deux modes de detection :
\begin{itemize}[leftmargin=*,nosep]
  \item \textbf{Directory listing} : liste le repertoire (simple, par defaut).
  \item \textbf{File notification} : s'abonne aux evenements du cloud (scalable, gros volumes).
\end{itemize}

\begin{lstlisting}[style=pystyle,caption={Auto Loader avec evolution de schema}]
df = (spark.readStream
        .format("cloudFiles")
        .option("cloudFiles.format", "json")
        .option("cloudFiles.schemaLocation", "/chk/schema")
        .load("s3://bucket/entrant/"))

(df.writeStream
   .option("checkpointLocation", "/chk/ventes")
   .option("mergeSchema", "true")
   .trigger(availableNow=True)   # mode batch : traite tout puis s'arrete
   .toTable("mon_catalogue.bronze.evenements"))
\end{lstlisting}

\subsection{COPY INTO vs Auto Loader : le choix cle}

\begin{center}
\renewcommand{\arraystretch}{1.3}
\begin{tabular}{@{}p{4.5cm}p{4.7cm}p{4.7cm}@{}}
\toprule
 & \textbf{COPY INTO} & \textbf{Auto Loader} \\
\midrule
Nature & SQL, batch & Structured Streaming \\
Volume de fichiers & Milliers (modere) & Millions (tres scalable) \\
Detection nouveaux fichiers & Suivi interne des fichiers charges & Directory listing ou file notification \\
Schema evolution & \texttt{mergeSchema} & Automatique + \texttt{schemaLocation} \\
Ideal pour & Chargements periodiques simples & Flux continu / gros volumes \\
\bottomrule
\end{tabular}
\end{center}

\subsection{Lakeflow Connect}

Connecteurs \textbf{managed} (sources SaaS/BD d'entreprise : Salesforce, SQL Server...) et
\textbf{standard}. Ingestion fiable et gouvernee dans des tables UC, y compris donnees
semi-structurees (JSON, nested).

\begin{piege}
\textbf{Auto Loader} = gros volumes, streaming, evolution de schema auto. \textbf{COPY INTO}
= batch SQL idempotent, volumes moderes. \textbf{Lakeflow Connect} = sources d'entreprise
managees. Ne les confonds pas : le choix se justifie par volume + frequence + type de source.
\end{piege}

\begin{piege}[: logs d'audit]
Les \textbf{audit logs} Databricks sont livres en \textbf{JSON}, generalement sous 15 min,
et de nouvelles livraisons \textbf{peuvent ecraser} des fichiers existants (ce n'est pas
CSV, pas Parquet, pas hebdomadaire).
\end{piege}

\begin{qcm}
Besoin : ingerer en continu des millions de petits fichiers JSON arrivant en permanence,
avec evolution de schema automatique. Quelle methode ?

\smallskip
A. COPY INTO en batch nocturne \hfill
B. \textbf{Auto Loader (cloudFiles) en streaming} \\
C. INSERT manuel par notebook \hfill
D. Lakeflow Connect vers Salesforce

\smallskip
\textit{Reponse : \textbf{B}. Auto Loader gere volume massif + schema evolution + checkpointing.}
\end{qcm}

% ============================================ S3 TRANSFORMATION
\section{Data Transformation and Modeling (22\%)}

\subsection{Architecture medaillon}

\begin{center}
\renewcommand{\arraystretch}{1.3}
\begin{tabular}{@{}p{2.5cm}p{12cm}@{}}
\toprule
\textbf{Couche} & \textbf{Role} \\
\midrule
Bronze & Donnees brutes ingerees telles quelles (append). \\
Silver & Donnees nettoyees : nulls traites, types standardises, dedupliquees, jointures. \\
Gold & Donnees agregees pretes pour le BI/analytics (tables, vues, vues materialisees). \\
\bottomrule
\end{tabular}
\end{center}

\subsection{Nettoyage Bronze $\rightarrow$ Silver}

\begin{lstlisting}[style=pystyle,caption={Lire bronze, nettoyer, ecrire silver}]
from pyspark.sql import functions as F

bronze = spark.read.table("cat.bronze.clients")

silver = (bronze
    .na.drop(subset=["client_id"])                 # supprimer nulls cles
    .withColumn("date", F.to_date("date_str"))     # standardiser le type
    .dropDuplicates(["client_id"]))                # deduplication

silver.write.mode("overwrite").saveAsTable("cat.silver.clients")
\end{lstlisting}

\subsection{Jointures (joins)}

\begin{itemize}[leftmargin=*,nosep]
  \item \texttt{inner}, \texttt{left}, \texttt{right}, \texttt{full}, \texttt{cross}.
  \item \textbf{Broadcast join} : diffuse la petite table sur tous les executors $\rightarrow$
  evite un shuffle couteux. Controle par \texttt{spark.sql.autoBroadcastJoinThreshold}.
  \item Jointure sur \textbf{cles multiples}, puis \texttt{union} / \texttt{unionAll}.
\end{itemize}

\begin{lstlisting}[style=pystyle,caption={Broadcast join et agregations}]
from pyspark.sql.functions import broadcast, approx_count_distinct, mean

df = grandes.join(broadcast(petites), "cle", "left")

agg = (df.groupBy("region")
         .agg(approx_count_distinct("client_id").alias("clients"),
              mean("montant").alias("panier_moyen")))
\end{lstlisting}

\subsection{Manipulations de colonnes/lignes}

Ajouter/supprimer/renommer colonnes, filtres, split, \textbf{exploser des tableaux}
(\texttt{explode}), deduplication, agregations (\texttt{count}, \texttt{approx\_count\_distinct},
\texttt{mean}, \texttt{summary}).

\subsection{Parametres de tuning a connaitre}

\begin{center}
\renewcommand{\arraystretch}{1.3}
\begin{tabular}{@{}p{6.5cm}p{8cm}@{}}
\toprule
\textbf{Parametre} & \textbf{Effet} \\
\midrule
\texttt{spark.sql.shuffle.partitions} & Nombre de partitions apres un shuffle (join/groupBy). \\
\texttt{spark.default.parallelism} & Parallelisme par defaut des RDD. \\
\texttt{spark.executor.memory} / \texttt{driver.memory} & Memoire allouee par executor / driver. \\
\texttt{spark.sql.autoBroadcastJoinThreshold} & Seuil sous lequel une table est broadcastee. \\
\bottomrule
\end{tabular}
\end{center}

\subsection{Objets de la couche Gold (UC)}

\begin{center}
\renewcommand{\arraystretch}{1.3}
\begin{tabular}{@{}p{4cm}p{10.5cm}@{}}
\toprule
\textbf{Objet} & \textbf{Description} \\
\midrule
Table (managed) & Donnees + metadonnees gerees par UC. \\
View & Requete sauvegardee, calculee a la lecture (pas de stockage). \\
Materialized view & Resultat precalcule et rafraichi $\rightarrow$ requetes BI rapides. \\
Streaming table & Table alimentee en continu par un flux (ingestion/transfo incrementale). \\
\bottomrule
\end{tabular}
\end{center}

\begin{piege}
\textbf{View} = recalculee a chaque lecture (toujours fraiche, plus lente).
\textbf{Materialized view} = precalculee (rapide, mais doit etre rafraichie).
\textbf{Streaming table} = mise a jour incrementale continue. Sais distinguer les trois usages.
\end{piege}

\begin{qcm}
Une equipe BI veut des tableaux de bord rapides sur des agregats couteux recalcules souvent.
Quel objet Gold ?

\smallskip
A. Vue standard \hfill
B. \textbf{Vue materialisee} \hfill
C. Table temporaire \hfill
D. CTE

\smallskip
\textit{Reponse : \textbf{B}. La vue materialisee precalcule et stocke le resultat $\rightarrow$
lectures BI rapides, contrairement a la vue standard recalculee a chaque fois.}
\end{qcm}

% ============================================ S4 LAKEFLOW JOBS
\section{Working with Lakeflow Jobs (16\%)}

\subsection{Orchestration par DAG}

Un \textbf{Lakeflow Job} (ex-Workflow) est un graphe oriente acyclique (DAG) de \textbf{taches}
avec dependances. Types de taches : \textbf{notebook}, \textbf{SQL query}, \textbf{dashboard},
\textbf{pipeline} (Lakeflow Spark Declarative Pipeline).

\subsection{Control flow}

\begin{itemize}[leftmargin=*,nosep]
  \item \textbf{Retries} : reessais automatiques en cas d'echec.
  \item \textbf{Conditional / branching} : tache \texttt{If/else} selon une condition.
  \item \textbf{Looping} : \texttt{For each} sur une liste de valeurs.
  \item \textbf{Dependances} : une tache s'execute apres la reussite de ses parents.
\end{itemize}

\subsection{Triggers (declencheurs)}

\begin{center}
\renewcommand{\arraystretch}{1.3}
\begin{tabular}{@{}p{4cm}p{10.5cm}@{}}
\toprule
\textbf{Trigger} & \textbf{Declenche quand...} \\
\midrule
Scheduled (cron) & Base sur le temps : a intervalle fixe (ex. toutes les heures). \\
File arrival & Un nouveau fichier arrive a un emplacement de stockage. \\
Table update & Une table amont est mise a jour (data-driven). \\
\bottomrule
\end{tabular}
\end{center}

\begin{piege}
\textbf{Time-based} (scheduled) : pour des donnees a horaire previsible.
\textbf{Data-driven} (file arrival / table update) : pour reagir a la \emph{disponibilite
reelle} des donnees, sans attendre une heure fixe ni interroger inutilement.
\end{piege}

\begin{qcm}
Les donnees sources arrivent a des heures imprevisibles. On veut lancer le pipeline des
qu'un fichier est depose, sans polling. Quel trigger ?

\smallskip
A. Scheduled toutes les 5 min \hfill
B. \textbf{File arrival} \hfill
C. Manuel \hfill
D. Table update sur une table sans lien

\smallskip
\textit{Reponse : \textbf{B}. Le trigger file arrival reagit a l'arrivee du fichier -- data-driven,
plus efficace qu'un cron serre.}
\end{qcm}

% ============================================ S5 CI/CD
\section{Implementing CI/CD (10\%)}

\subsection{Git Folders (ex-Repos)}

Dans le workspace : creer/changer de branche, commit, push, ouvrir une pull request via
l'integration Git. Workflow de dev versionne.

\subsection{Automation Bundles (ex-Databricks Asset Bundles / DAB)}

Format declaratif (\texttt{databricks.yml}) qui \textbf{package et promeut} du code + des
ressources (Lakeflow Jobs, pipelines declaratifs, etc.) a travers \textbf{dev / test / prod}.
Gestion des differences d'environnement par \textbf{variables et overrides} (targets).

\begin{lstlisting}[style=sqlstyle,caption={Structure d'un bundle (databricks.yml, simplifie)}]
bundle:
  name: mon_pipeline_etl

targets:
  dev:
    default: true
    workspace: { host: https://dev.cloud.databricks.com }
  prod:
    workspace: { host: https://prod.cloud.databricks.com }
    variables: { catalogue: prod_cat }
\end{lstlisting}

\subsection{Databricks CLI}

Valide, deploie et gere les bundles dans des workflows CI/CD automatises :
\begin{itemize}[leftmargin=*,nosep]
  \item \texttt{databricks bundle validate} -- verifie la config.
  \item \texttt{databricks bundle deploy -t prod} -- deploie sur la cible.
  \item \texttt{databricks bundle run} -- lance un job du bundle.
\end{itemize}

\begin{piege}
Pour \textbf{deployer, versionner et orchestrer} des pipelines ETL de facon modulaire et
repetable (CI/CD) : la bonne reponse est \textbf{Automation Bundles + Git + CLI}, \emph{pas}
des wheels dans des Volumes, ni la revision history des notebooks comme systeme de versioning.
\end{piege}

\begin{qcm}
Une equipe veut un moyen modulaire de deployer, versionner et orchestrer ses pipelines ETL
avec CI/CD et repetabilite. Quelle fonctionnalite ?

\smallskip
A. Modeles UC pour representer les jobs ETL \hfill
B. Wheels dans des Volumes UC \\
C. Notebook monte sur Volume + Jobs API v2 \hfill
D. \textbf{Automation Bundles (DABs) versionnes dans Git, promus par CI/CD}

\smallskip
\textit{Reponse : \textbf{D}. Les bundles definissent ressources + code, versionnes dans Git,
promus dev$\to$test$\to$prod par CI/CD automatise.}
\end{qcm}

% ============================================ S6 TROUBLESHOOTING
\section{Troubleshooting, Monitoring, Optimization (10\%)}

\subsection{Monitoring via l'UI Lakeflow Jobs}

\begin{itemize}[leftmargin=*,nosep]
  \item \textbf{Run history} : comparer les temps d'execution actuels aux baselines historiques
  (detecter les tendances de degradation).
  \item Interpreter les statuts de jobs, lire le \textbf{DAG} pour reperer les blocages amont.
  \item Suivre les temps d'execution et les taux d'echec.
\end{itemize}

\subsection{Bottlenecks dans la Spark UI}

\begin{center}
\renewcommand{\arraystretch}{1.3}
\begin{tabular}{@{}p{3.5cm}p{11cm}@{}}
\toprule
\textbf{Probleme} & \textbf{Symptome / metrique} \\
\midrule
Data skew & Une tache beaucoup plus longue que les autres ; Max shuffle read $\gg$ Median. \\
Shuffling & Beaucoup d'echange de donnees entre executors (joins/groupBy). \\
Disk spilling & Donnees ecrites sur disque faute de memoire (spill memory/disk). \\
\bottomrule
\end{tabular}
\end{center}

\begin{piege}[: data skew]
Symptome typique : \emph{la plupart des taches finissent en <30s mais une prend >10 min,
Max shuffle read $\gg$ Median}. C'est du \textbf{data skew}. La solution robuste :
verifier que l'\textbf{Adaptive Query Execution (AQE) avec skew join handling} est actif
(il decoupe la partition surdimensionnee au runtime). Ajouter des executors ne corrige
pas le desequilibre.
\end{piege}

\subsection{Optimisations Delta}

\begin{itemize}[leftmargin=*,nosep]
  \item \textbf{Liquid Clustering} : remplace le partitionnement + ZORDER, s'adapte
  automatiquement aux patterns de requetes (\texttt{CLUSTER BY}).
  \item \textbf{Predictive Optimization} : Databricks execute automatiquement \texttt{OPTIMIZE}
  et \texttt{VACUUM} au bon moment, sans reglage manuel.
\end{itemize}

\subsection{Diagnostics cluster}

Savoir diagnostiquer : echecs de demarrage de cluster, conflits de librairies,
erreurs out-of-memory (OOM).

\begin{qcm}
Dans la Spark UI, une tache prend 10 min alors que les autres finissent en 30s, avec un
Max shuffle read a 5 GB contre 400 MB en median. Que faire ?

\smallskip
A. Agrandir le cluster \hfill
B. \textbf{Confirmer que l'AQE avec skew join handling est actif} \\
C. Reduire \texttt{shuffle.partitions} \hfill
D. Repartitionner manuellement avec une salt key

\smallskip
\textit{Reponse : \textbf{B}. C'est du data skew ; l'AQE decoupe automatiquement la partition
surdimensionnee. (D peut aider mais B est la solution native recommandee.)}
\end{qcm}

% ============================================ S7 GOVERNANCE
\section{Governance and Security (15\%)}

\subsection{Tables managed vs external}

\begin{center}
\renewcommand{\arraystretch}{1.3}
\begin{tabular}{@{}p{3cm}p{5.5cm}p{5.5cm}@{}}
\toprule
 & \textbf{Managed} & \textbf{External} \\
\midrule
Donnees & Gerees par UC (emplacement gere) & A un emplacement externe que tu specifies \\
DROP TABLE & Supprime donnees + metadonnees & Supprime seulement les metadonnees \\
Usage & Cas par defaut, cycle de vie simple & Donnees partagees / possedees ailleurs \\
\bottomrule
\end{tabular}
\end{center}

\begin{lstlisting}[style=sqlstyle,caption={Creation et conversion de tables}]
-- Table managed
CREATE TABLE cat.silver.clients (id INT, nom STRING);

-- Table external (LOCATION explicite)
CREATE TABLE cat.silver.ext_clients (id INT)
LOCATION 's3://bucket/clients/';
\end{lstlisting}

\subsection{Controle d'acces : GRANT / REVOKE / DENY}

Appliques a des \textbf{principals} (users, groups, service principals) aux bons niveaux de
la hierarchie (catalog / schema / table).

\begin{lstlisting}[style=sqlstyle,caption={Privileges Unity Catalog}]
GRANT SELECT ON TABLE cat.silver.clients TO `analystes`;
GRANT USAGE  ON SCHEMA cat.silver          TO `analystes`;
REVOKE SELECT ON TABLE cat.silver.clients FROM `stagiaires`;
DENY   MODIFY ON TABLE cat.silver.clients TO `lecture_seule`;
\end{lstlisting}

\subsection{Masquage et securite au niveau ligne/colonne}

\begin{itemize}[leftmargin=*,nosep]
  \item \textbf{Column masking} : masquer une colonne sensible selon le groupe de l'utilisateur.
  \item \textbf{Row-level security (RLS)} : filtrer les lignes visibles selon l'utilisateur.
  \item \textbf{ABAC policies} (Attribute-Based Access Control) : controler centralement le
  filtrage de lignes et le masquage de colonnes pour les donnees sensibles, via des attributs.
\end{itemize}

\begin{piege}
\textbf{Column masking} agit sur les \emph{colonnes} (cache une valeur).
\textbf{Row-level security} agit sur les \emph{lignes} (cache des enregistrements).
\textbf{ABAC} = gestion centralisee des deux via des attributs/tags, a l'echelle.
\end{piege}

\begin{qcm}
On veut que les analystes voient la table clients mais que le numero de carte bancaire soit
masque pour les groupes non autorises. Quel mecanisme ?

\smallskip
A. Row-level security \hfill
B. \textbf{Column-level masking} \hfill
C. REVOKE sur la table \hfill
D. Table external

\smallskip
\textit{Reponse : \textbf{B}. Le masquage de colonne cache une valeur precise selon le groupe,
sans retirer l'acces a toute la table.}
\end{qcm}

% ============================================ HANDS-ON
\section{Checklist hands-on (a savoir faire les yeux fermes)}

Le guide de preparation insiste : \textbf{lire ne suffit pas}. Construis chacune de ces taches
dans Databricks Free Edition (< 30 min chacune). Toute tache que tu ne sais pas faire = ta
prochaine cible de revision.

\begin{center}
\renewcommand{\arraystretch}{1.4}
\begin{longtable}{@{}p{0.6cm}p{10.4cm}p{3cm}@{}}
\toprule
 & \textbf{Tache} & \textbf{Objectif} \\
\midrule
\endhead
$\square$ & Ingerer des fichiers d'un dossier cloud avec \textbf{Auto Loader} (schema evolution) & \S2 \\
$\square$ & Charger des CSV avec \textbf{COPY INTO} dans une table UC & \S2 \\
$\square$ & Nettoyer bronze $\rightarrow$ silver (nulls, types, \textbf{dropDuplicates}) & \S3 \\
$\square$ & Faire un \textbf{broadcast join} + agregations (\texttt{approx\_count\_distinct}) & \S3 \\
$\square$ & Creer une \textbf{materialized view} et une \textbf{streaming table} & \S3 \\
$\square$ & Construire un \textbf{Lakeflow Job} multi-taches avec dependances + retries & \S4 \\
$\square$ & Configurer un trigger \textbf{file arrival} et un trigger \textbf{scheduled} & \S4 \\
$\square$ & Deployer un \textbf{Automation Bundle} sur une cible via la \textbf{CLI} & \S5 \\
$\square$ & Lire la \textbf{Spark UI} et identifier un \textbf{data skew} & \S6 \\
$\square$ & Appliquer \textbf{GRANT / REVOKE} et un \textbf{column mask} en SQL & \S7 \\
$\square$ & Creer une table \textbf{managed} et une \textbf{external}, puis les DROP (observer la difference) & \S7 \\
\bottomrule
\end{longtable}
\end{center}

% ============================================ QUESTIONS OFFICIELLES
\section{Les 5 questions officielles du guide (corrigees)}

Ces questions proviennent du guide officiel. Comprends \emph{pourquoi} chaque reponse est correcte.

\begin{enumerate}[leftmargin=*]
  \item \textbf{Data skew (Spark UI).} Tache 10 min vs 30s, Max shuffle read $\gg$ median.
        \textbf{Reponse B} : AQE + skew join handling decoupe la partition au runtime.
  \item \textbf{Compute / plateforme.} Iteration rapide, rollback, audit, source unique IA+BI.
        \textbf{Reponse B} : Delta Lake (ACID + time travel) gouverne par Unity Catalog.
  \item \textbf{Audit logs.} Format et comportement des logs d'audit Databricks.
        \textbf{Reponse A} : JSON, < 15 min, les livraisons peuvent ecraser des fichiers.
  \item \textbf{Cluster pour analystes SQL ad hoc.} Concurrence + demarrage rapide + cout.
        \textbf{Reponse C} : cluster high-concurrency avec autoscaling.
  \item \textbf{CI/CD pipelines ETL modulaires.} Deployer, versionner, orchestrer.
        \textbf{Reponse D} : Automation Bundles (DABs), versionnes dans Git, promus par CI/CD.
\end{enumerate}

% ============================================ MEMO
\section*{Annexe : memo derniere minute (a relire 1h avant)}
\addcontentsline{toc}{section}{Annexe : memo derniere minute}

\begin{cle}[: les reflexes qui font gagner des points]
\begin{itemize}[leftmargin=*,nosep]
  \item \textbf{Noms actuels} : Lakeflow Jobs, Lakeflow Spark Declarative Pipelines, Git Folders,
  Automation Bundles, Lakeflow Connect. Jamais Workflows / DLT / Repos / DAB dans une bonne reponse.
  \item \textbf{Auto Loader} = streaming, gros volumes, schema evolution auto. \textbf{COPY INTO}
  = batch SQL idempotent. \textbf{Lakeflow Connect} = sources d'entreprise managees.
  \item \textbf{Delta} = ACID + time travel (rollback/audit). \textbf{Unity Catalog} = gouvernance,
  hierarchie catalog $\to$ schema $\to$ objet.
  \item \textbf{Medaillon} : bronze (brut) $\to$ silver (propre) $\to$ gold (agrege BI).
  \item \textbf{View} recalculee / \textbf{Materialized view} precalculee / \textbf{Streaming table} incrementale.
  \item \textbf{Data skew} $\rightarrow$ AQE + skew join handling.
  \item \textbf{Liquid Clustering} (remplace partition+ZORDER) ; \textbf{Predictive Optimization}
  (OPTIMIZE/VACUUM auto).
  \item \textbf{Triggers} : scheduled (temps) vs file arrival / table update (data-driven).
  \item \textbf{Bundles + Git + CLI} = CI/CD. \texttt{validate} $\to$ \texttt{deploy} $\to$ \texttt{run}.
  \item \textbf{Managed} : DROP supprime les donnees. \textbf{External} : DROP garde les donnees.
  \item \textbf{GRANT/REVOKE/DENY} sur principals ; \textbf{column mask} (colonne) / \textbf{RLS}
  (ligne) / \textbf{ABAC} (centralise par attributs).
\end{itemize}
\end{cle}

\vfill
\begin{center}
\textit{Verifie toujours sur} \url{https://docs.databricks.com} \textit{-- bonne chance !}
\end{center}

\end{document}
